\section{Challenges with AI-agent analysis}
\label{sec:challenges}

This section is maintained as a living record of unresolved or recurring issues encountered during the analysis. Because this document is an internal note rather than a final publication, we prioritize explicit tracking of weaknesses, mitigation status, and open analysis challenges. Practical lessons from AI--physicist collaboration are discussed in Sec.~\ref{sec:context}.

\subsection{Current limitations of the thrust-only baseline}

The present release is robust enough for a preliminary thrust result, but several limitations remain:
\begin{itemize}
\item for events with exactly 3 selected particles, the thrust algorithm selects the highest-magnitude particle as the axis rather than enumerating all possible hemisphere assignments; this approximation is justified because (i) 3-particle hadronic events are rare in $Z \to q\bar{q}$ decays (impact on thrust value typically $\lesssim 0.1\%$ for $T > 0.5$), and (ii) the response-matrix migration structure dominates uncertainty in the endpoint bins where such events concentrate;
\item hadronic-event correction uncertainty is propagated as an independent component ($\Delta_{\mathrm{had}}$, central-range mean $\approx 3.2\%$); the endpoint bins remain sensitive to correction-factor extrapolation;
\item systematic inter-bin correlations use the statistical correlation matrix as a template ($\mathrm{Cov}^{(c)}_{ij} = \Delta_c(i)\,\Delta_c(j)\,\rho^{\mathrm{stat}}_{ij}$); this is an approximation that assumes systematic and statistical correlations share the same bin-to-bin structure;
\item endpoint thrust bins remain statistically and modeling sensitive;
\item analysis scope is currently one-dimensional (thrust only), while the future target includes the thrust-family observables and potentially multidimensional unfolding.
\end{itemize}

\subsection{Persistent analysis challenges}

Several recurrent issues have been tracked through the development of this analysis. Plot-style consistency across scripts has been an ongoing concern, with centralized style enforcement in producer scripts now largely resolving the issue. Method-to-code traceability gaps are addressed by the workflow and provenance tables in Sec.~\ref{sec:method} and Sec.~\ref{sec:results}, with \path{workflow_summary.json} serving as the canonical numeric record. The most significant open challenge is endpoint-bin instability: the final bin ($0.99<T<1.00$) is sensitive to MC modeling, hadronic correction extrapolation, and normalization convention differences, and dominates the maximum per-bin uncertainty. This is mitigated by reporting central-window metrics ($0.60\le T<0.95$) alongside full-range figures throughout this paper.

\subsection{Update protocol for this section}

For each major workflow iteration, this section should record:
\begin{itemize}
\item newly discovered persistent issues;
\item whether previous issues moved from Open to In progress/Resolved;
\item concrete next action and code location for unresolved items.
\end{itemize}

Keeping this log in the main body is intentional: it makes assumptions and remaining risks visible to internal readers and future maintainers.

\subsection{Scaling human oversight: toward systematic validation}
\label{sec:scaling_oversight}

A recurring theme in this work is that the set of things that can go wrong silently grows rapidly with analysis complexity. The one-dimensional thrust measurement already required validation of event selection, unfolding closure, systematic propagation, covariance-matrix properties, and figure/table accuracy. Extending the program to additional observables, multidimensional unfolding, or theory--experiment feedback loops will multiply these failure modes.

The current mitigation relies on structured sign-off points, automated integrity checks (file presence, covariance positive semi-definiteness, fit-mask validity), and version-controlled provenance. While adequate for a single-observable proof of concept, this approach does not scale without further investment in systematic testing infrastructure.

Software engineering and machine-learning operations (MLOps) communities have partially addressed analogous challenges through continuous-integration and continuous-deployment (CI/CD) pipelines, where automated test suites run on every code change and failures block deployment. Adapting these patterns to physics analysis workflows is a natural next step: for example, automated regression tests could verify that modifying one systematic variation does not silently shift the unfolded spectrum outside a predefined tolerance. Physics-specific assertions --- such as verifying that the unfolded distribution integrates to unity, that the response matrix is properly normalized, or that refolded spectra agree with detector-level data --- could be encoded as unit tests that run after each agent iteration.

The key unsolved challenge is encoding implicit physics priors as testable assertions. Many of the validation checks performed by the physicist in this analysis (Sec.~\ref{sec:ai_agent}) rely on domain intuition that is difficult to formalize: recognizing that a $\varphi$ excess indicates a detector artifact, judging whether a per-bin uncertainty is ``physically reasonable,'' or deciding that a $23\%$ shift constitutes an observable redefinition rather than a systematic effect. Developing methods to translate such expert knowledge into machine-executable tests --- whether through formal specification languages, learned anomaly detectors, or structured physicist--agent dialogue protocols --- is an important direction for making AI-assisted experimental physics robust and scalable.
